\section{Background}
\label{sec:background}

This first section briefly covers [\dots]. The following sections introduce [\dots], and finally how [\dots].

\subsection{things in background}
\label{sec:thing1}


\lipsum[1]


\begin{figure}[h]
	\centering
	\includegraphics[width=0.3\textwidth]{resources/temp.png}
	\caption{A simple graphic}
	\label{fig:simple}	
\end{figure}


\subsubsection*{Non-numbered subsubsection}


\begin{description}
	\item [First] This is the first Item of a description list, could be useful if you need to name each procedure in a system, or in other ways list various components, while giving a short description of them.

	\item [Second] item in a list.
	

\end{description}

\lipsum[5]

\pagebreak

\todo{Remember that its easier to segment a paper into various sub-components, are you inserting large tables, tikz-pictures or other, then it might be wise to place them in a separeate tex and link it instead of having a gigantic single text file.}

\begin{figure}[htb]
	\centering
	\resizebox{\textwidth}{!}{\tikzset{
	data/.style={
		draw,
		rectangle split,
		rectangle split parts=2,
		text centered,
	},
	data+/.style={
		data,
		rectangle split every empty part={},% resets empty-part macro (explanation below)
		rectangle split empty part width=\widthof{#1},
		rectangle split empty part height=\heightof{#1},
		rectangle split empty part depth=\depthof{#1},
	},
}


\begin{tikzpicture}[background rectangle/.style={fill=green!15, rounded
	corners, draw=green!50}, show background rectangle, split/.style={rectangle split, rectangle split horizontal, rectangle split parts=#1, draw},splitv/.style={rectangle split, rectangle split parts=#1, draw}]



\node (R) {Root};
\node[split=1, above=0.25cm of R, fill=black!60] (A) {};


\node [data, right=1.0cm of A, fill=black!30] (B) { \nodepart{second} marked = True};
\draw[-stealth] ([yshift=0.5cm]A) -- (B);
\node [data, right=1.0cm of B, fill=black!30] (C) { \nodepart{second} marked = True};
\draw[-stealth, transform canvas={yshift=0.25cm}] (B) -- (C);
\node [data, right=1.0cm of C, fill=black!30] (D) { \nodepart{second} marked = True};
\draw[-stealth, transform canvas={yshift=0.25cm}] (C) -- (D);
\draw[-stealth, transform canvas={yshift=0.25cm}] (D) to[out=5, in=20] (B);

\node [data, below=1cm of B, fill=white] (B1) { \nodepart{second} marked = False};
\node [data, right=1.0cm of B1, fill=white] (C1) { \nodepart{second} marked = False};
\draw[-stealth, transform canvas={yshift=0.25cm}] (B1) -- (C1);
\node [data, right=1.0cm of C1, fill=black!30] (D1) { \nodepart{second} marked = True};
\draw[-stealth, transform canvas={yshift=0.25cm}] (C1) -- (D1);

\draw[-stealth, transform canvas={yshift=0.25cm}] (B1) to[out=20, in=-170,distance=0.5cm] (C);
\draw[-stealth, transform canvas={yshift=0.25cm}] (C) to[out=0, in=-190,distance=0.5cm] (D1);
\node[below=0.6cm of C1] {Illustration after a mark.};



\end{tikzpicture}}
	\caption{This image is a tex-document written in tikz.}
	\label{fig:tikz}
\end{figure}

\lipsum[2-4]

\pagebreak